\begin{solapartat}
\begin{minipage}[t]{0.5\linewidth}
\vspace{0pt}
\figura[0.95]{sol_fig1.png}
\end{minipage}\hfill
\begin{minipage}[t]{0.46\linewidth}
\vspace{0pt}
El vector camp magnètic és tangent a la línia de camp (direcció x). \punts{0,5}
\end{minipage}

\medskip
% A l'original, la puntuació d'aquesta frase és entre parèntesis: «(0.2 p)».
El sentit de la força que actua sobre el fil determina el sentit de gir. \punts{0,2}

El sentit de la força magnètica és perpendicular al pla format pel camp magnètic i la intensitat (cal avaluar el producte vectorial entre la intensitat i el camp magnètic). Segons el sistema de referència adjunt, la força «surt del pla del paper» (direcció z). \punts{0,3}
\end{solapartat}

\begin{solapartat}
\punts{0,5} $\vec{F} = I\left(\vec{l}\wedge\vec{B}\right) \Rightarrow F = IlB\sin 90° = 10\cdot 0,01\cdot 0,1 = 0,01\,N$\quad \punts{0,5}
\end{solapartat}
